\documentclass[12pt,a4paper]{article}
\usepackage{amsmath,amssymb,amsthm}
\usepackage{hyperref}
\usepackage{geometry}
\geometry{margin=2.5cm}
\usepackage{enumitem}
\usepackage{booktabs}

\newtheorem{theorem}{Theorem}[section]
\newtheorem{lemma}[theorem]{Lemma}
\newtheorem{corollary}[theorem]{Corollary}
\newtheorem{proposition}[theorem]{Proposition}
\theoremstyle{definition}
\newtheorem{definition}[theorem]{Definition}
\theoremstyle{remark}
\newtheorem{remark}[theorem]{Remark}

\title{The Measurement Problem Dissolved:\\
Recognition IS Measurement in Recognition Science}

\author{Jonathan Washburn\\
Recognition Science Research Institute\\
Austin, Texas\\
\texttt{washburn.jonathan@gmail.com}}

\date{March 2026}

\begin{document}
\maketitle

\begin{abstract}
We dissolve the quantum measurement problem within Recognition
Science (RS).  The problem---how and when does the wave function
collapse?---arises from treating measurement as an external process
applied to a quantum system.  In RS, measurement \emph{is}
recognition: the discrete ledger commits to a definite outcome at
each 8-tick epoch boundary.  The ledger has two phases:
\emph{uncommitted} (superposition of branches with complex
amplitudes) and \emph{committed} (single definite outcome).  The
transition is the irreversible ledger commit operation---this is
``collapse.''  The Born rule $P(i) = |c_i|^2$ follows from five
constraints: non-negativity, normalization, $J$-cost phase
invariance, additivity under coarse-graining, and continuity.
We show that $P(i) = |c_i|^2$ is the unique
probability assignment satisfying all five.  Definite outcomes are
forced by the discreteness of the ledger (one entry per tick).
The 8-tick epoch ($2^D = 8$ for $D = 3$) provides the minimal
complete measurement cycle.  Schr\"odinger's cat,
Wigner's friend, and the quantum-to-classical transition are all
resolved by the same mechanism.
\end{abstract}

%% ============================================================
\section{Introduction}
\label{sec:intro}
%% ============================================================

The measurement problem, identified by von Neumann in
1932~\cite{vN}, remains the central interpretive puzzle of quantum
mechanics.  Every interpretation must answer three questions:
\begin{enumerate}[nosep]
  \item \textbf{When} does the wave function collapse?
  \item \textbf{How} does a superposition become a definite
  outcome?
  \item \textbf{Why} is the probability rule $P = |\psi|^2$?
\end{enumerate}

The standard formalism~\cite{Bell} contains two evolution laws:
unitary Schr\"odinger evolution (deterministic, reversible, linear)
and projection (stochastic, irreversible, nonlinear).  No mechanism
specifies when one law replaces the other.

Proposed solutions include:
\begin{itemize}[nosep]
  \item \textbf{Copenhagen}: Collapse occurs upon ``observation''
  (undefined).
  \item \textbf{Many-Worlds}~\cite{Everett}: No collapse;
  all branches exist.
  \item \textbf{Bohmian mechanics}~\cite{Bohm}: Deterministic
  particles guided by $\psi$; collapse is epistemic.
  \item \textbf{GRW}~\cite{GRW}: Spontaneous collapse at random
  times.
  \item \textbf{QBism}~\cite{Fuchs}: $\psi$ encodes agent beliefs.
\end{itemize}

Recognition Science~\cite{RS_Axioms} does not \emph{solve} the
measurement problem within the existing framework; it \emph{dissolves}
it by providing a framework in which the problem does not arise.

%% ============================================================
\section{The Cost Functional}
\label{sec:cost}
%% ============================================================

\begin{definition}[RS Cost Functional]
For $x > 0$:
\begin{equation}
  J(x) = \tfrac{1}{2}(x + x^{-1}) - 1.
\end{equation}
\end{definition}

\noindent The functional $J$ is the unique continuous solution to
three axioms:\\
\textbf{Axiom A1 (Normalization):} $J(1) = 0$.\\
\textbf{Axiom A2 (Recognition Composition Law, RCL):}
$J(xy) + J(x/y) = 2J(x)J(y) + 2J(x) + 2J(y)$ for all $x, y > 0$.\\
\textbf{Axiom A3 (Calibration):} $J''_{\log}(0) = 1$, where
$J_{\log}(t) := J(e^t)$.

\begin{theorem}[Cost Uniqueness]
\label{thm:unique}
The unique continuous solution to \textup{(A1)--(A3)} is
$J(x) = \tfrac{1}{2}(x + x^{-1}) - 1$.
\end{theorem}

\begin{proof}[Proof sketch]
Setting $f(t) = J_{\log}(t) + 1$ and rewriting \textup{(A2)} gives the
d'Alembert equation $f(s+t) + f(s-t) = 2f(s)f(t)$.  By the Acz\'el
classification theorem~\cite{Aczel1966}, the unique continuous solution
with $f(0) = 1$ and $f''(0) = 1$ is $f(t) = \cosh t$.
\end{proof}

\begin{lemma}[Phase Invariance]
\label{lem:phase}
$J$ depends only on the modulus: for $z = r\,e^{i\theta}$ with
$r > 0$,
\begin{equation}
  J(|z|) = J(r) \quad \text{(independent of $\theta$)}.
\end{equation}
\end{lemma}

%% ============================================================
\section{The Ledger Commit Model}
\label{sec:commit}
%% ============================================================

\begin{definition}[Uncommitted Ledger State]
An uncommitted ledger state is a superposition of branches:
\begin{equation}
  |\Psi\rangle = \sum_i c_i\,|b_i\rangle,
  \quad c_i \in \mathbb{C},
  \quad \sum_i |c_i|^2 = 1,
\end{equation}
where $|b_i\rangle$ are orthogonal branch states with complex
amplitudes $c_i$.
\end{definition}

\begin{definition}[Committed Ledger State]
A committed ledger state is a single definite outcome
$|b_k\rangle$.  All bond multipliers take integer values---no
superposition.
\end{definition}

\begin{definition}[Collapse = Commit]
The transition
$\text{uncommitted} \to \text{committed}$ is the ledger commit
operation.  At every 8-tick epoch boundary ($t \bmod 8 = 0$), the
ledger commits to one branch $|b_k\rangle$ with probability $P(k)$
determined by the amplitudes $\{c_i\}$.  What $P(k)$ is---and the
fact that it equals $|c_k|^2$---follows from the uniqueness theorem
in Section~\ref{sec:born}; it is not part of this definition.
\end{definition}

\begin{proposition}[Measurement Irreversibility --- Structural Argument]
\label{prop:irreversible}
The commit operation is irreversible: a committed ledger entry
cannot be uncommitted.
\end{proposition}

\begin{proof}
The ledger's double-entry bookkeeping records every commit as a
balanced transaction (debit and credit).  Reversing a commit would
require creating an unbalanced entry (credit without debit).  An
unbalanced entry violates the Recognition Composition Law, giving
$J \to \infty$.  Therefore, the commit is permanent.
\end{proof}

\begin{remark}
Irreversibility resolves the puzzle of ``von Neumann's Process~1
vs.\ Process~2.''  Unitary evolution (Process~2) is the inter-tick
dynamics of the uncommitted ledger.  The commit (Process~1) occurs
at epoch boundaries.  Both are part of the same physical system;
they operate at different phases of the 8-tick cycle.
\end{remark}

%% ============================================================
\section{The Born Rule Derived}
\label{sec:born}
%% ============================================================

\begin{theorem}[Born Rule from $J$-Cost]
\label{thm:born}
Let $P: \mathbb{C}^n \to [0,1]^n$ assign probabilities to the
branches of an uncommitted ledger state
$|\Psi\rangle = \sum_i c_i |b_i\rangle$.  Suppose $P$ satisfies:
\begin{enumerate}[label=(\roman*),nosep]
  \item \textbf{Non-negativity}: $P(i) \geq 0$ for all $i$.
  \item \textbf{Normalization}: $\sum_i P(i) = 1$.
  \item \textbf{Phase invariance}: $P(i)$ depends on $c_i$
  only through $|c_i|$ ($J$ is phase-independent,
  Lemma~\ref{lem:phase}).
  \item \textbf{Additivity}: $P$ respects coarse-graining:
  grouping $k$ equal-amplitude branches into one outcome gives
  probability $k \cdot P(\text{single branch})$.
  \item \textbf{Continuity}: $P(i)$ is continuous in $|c_i|$.
\end{enumerate}
Then $P(i) = |c_i|^2$.
\end{theorem}

\begin{proof}
By (iii), $P(i) = f(|c_i|)$ for some function
$f:[0,\infty) \to [0,1]$.

\textit{Step 1 (equal amplitudes).}
Consider a state with $n$ equal-amplitude branches:
$c_i = e^{i\theta_i}/\sqrt{n}$.  By (iii), all $P(i)$ are equal.
By (ii), $nf(1/\sqrt{n}) = 1$, so
$f(1/\sqrt{n}) = 1/n = (1/\sqrt{n})^2$.

\textit{Step 2 (coarse-graining).}
Take $n$ branches of equal amplitude $1/\sqrt{n}$.  Group the
first $k$ branches into a single coarse-grained outcome.  By
additivity, the probability of the coarse-grained outcome is
$kf(1/\sqrt{n}) = k/n$.  By relabeling, this is also
$f(\sqrt{k/n})$ (since the coarse-grained amplitude is
$\sqrt{k/n}$).  Therefore
$f(\sqrt{k/n}) = k/n = (\sqrt{k/n})^2$ for all
integers $0 \leq k \leq n$.

\textit{Step 3 (density and continuity).}
The set $\{\sqrt{k/n} : k, n \in \mathbb{N},\; 0 \leq k \leq n\}$
is dense in $[0,1]$ (since $\{k/n\}$ is the set of rationals and
$\sqrt{\cdot}$ preserves density).  By condition (v) (continuity),
$f$ agrees with $r^2$ on a dense set and is continuous, so
$f(r) = r^2$ for all $r \in [0,1]$.  Therefore
$P(i) = |c_i|^2$.
\end{proof}

\begin{remark}[Interference from the Born Rule]
For a two-branch superposition $|\Psi\rangle = c_1|b_1\rangle +
c_2|b_2\rangle$, the Born rule gives the probability of measuring
the superposed amplitude $c_1 + c_2$ as:
\begin{equation}
  |c_1 + c_2|^2 = |c_1|^2 + |c_2|^2
  + 2\,\mathrm{Re}(c_1^*\,c_2).
\end{equation}
The cross-term $2\,\mathrm{Re}(c_1^*\,c_2)$ is the interference
term, responsible for quantum interference phenomena.  In the ledger,
this cross-term represents the cost of maintaining phase coherence
between two branches; it vanishes upon decoherence (commit).
\end{remark}

\begin{proposition}[Post-Measurement Definiteness]
\label{thm:definite}
After the commit, the ledger state is definite:
$P(\text{outcome}) = 1$.  No further branching occurs until the
next 8-tick epoch.
\end{proposition}

\begin{proof}[Structural argument]
By Definition~[Committed Ledger State], a committed state is a
single branch $|b_k\rangle$.  The probability assigned to $|b_k\rangle$
by the Born rule (Theorem~\ref{thm:born}) is $|1|^2 = 1$.
The definition of ``commit'' includes the commitment holding until
the next 8-tick epoch boundary; there is no second commit within the
same epoch.
\end{proof}

%% ============================================================
\section{Definite Outcomes from Discreteness}
\label{sec:discrete}
%% ============================================================

\begin{proposition}[No Superposition of Committed Entries]
\label{prop:nosuperposition}
Each committed ledger entry is a single definite integer-valued record.
Superposition describes the uncommitted inter-tick phase only.
\end{proposition}

\begin{proof}
The ledger is discrete: each cell contains exactly one
integer-valued entry.  An entry cannot be in a superposition of
values, just as a database record cannot be ``half 0 and half 1.''
The discreteness forces a definite outcome at every commit.
\end{proof}

\begin{remark}
This is the RS answer to ``How does a superposition become a
definite outcome?''  The answer: the superposition was never
``real'' in the ontic sense.  It describes the inter-tick dynamics
of an uncommitted state.  The committed state---the physically
real record---is always definite.
\end{remark}

%% ============================================================
\section{The 8-Tick Measurement Cycle}
\label{sec:8tick}
%% ============================================================

\begin{proposition}[Complete Measurement Cycle]
\label{prop:8tick}
One complete measurement cycle is one 8-tick epoch
($2^D = 8$ ticks for $D = 3$).
\end{proposition}

\begin{proof}
The $D$-cube has $2^D = 8$ vertices.  The phase structure
traverses $e^{2\pi i k/8}$ for $k = 0, 1, \ldots, 7$.  At the
half-period ($k = 4$): $e^{i\pi} = -1$, generating the complex
structure.  At the full period ($k = 8$): $e^{2\pi i} = 1$,
restoring real values.  The commit occurs at this boundary, when
the phase cycle returns to $+1$.
\end{proof}

\begin{corollary}
The measurement timescale is $\tau_{\mathrm{meas}} = 8\tau_0$,
where $\tau_0$ is the fundamental tick time.  At macroscopic scales,
this is negligibly short.
\end{corollary}

%% ============================================================
\section{Schr\"odinger's Cat Dissolved}
\label{sec:cat}
%% ============================================================

\begin{proposition}[No Macroscopic Superpositions]
\label{prop:cat}
Schr\"odinger's cat is never in a superposition of alive and dead.
The radioactive decay event commits at a specific 8-tick boundary,
and the cat's state follows deterministically from that commit.
\end{proposition}

\begin{proof}
The radioactive nucleus has an uncommitted state during each
8-tick epoch.  At the epoch boundary, the ledger commits:
either ``decayed'' or ``not decayed.''  This commit propagates
causally through the apparatus to the cat within subsequent ticks.
At no point does the macroscopic system enter a superposition.
The apparent superposition in the standard formalism arises from
extending unitary evolution past the commit boundary---a domain
error in the semiclassical approximation.
\end{proof}

%% ============================================================
\section{Wigner's Friend Resolved}
\label{sec:wigner}
%% ============================================================

\begin{proposition}[Wigner's Friend]
\label{prop:wigner}
The Wigner's friend paradox does not arise in RS.  Both Wigner and
his friend observe the same committed ledger entry.
\end{proposition}

\begin{proof}
The friend performs a measurement (commit) at epoch boundary $t_1$.
The result is recorded in the ledger as a definite entry.  When
Wigner later queries the friend, he reads the same committed
entry.  There is no ``superposition of the friend's knowledge'':
the ledger is a shared, objective record.  Wigner's ignorance
of the result before querying is epistemic, not ontic.
\end{proof}

\begin{remark}[Extended Wigner's Friend]
The same mechanism resolves the extended Wigner's friend
scenario of Frauchiger and Renner~\cite{FrauchigerRenner}.
Their paradox arises when two observers reason about each
other's quantum measurements and arrive at contradictory
conclusions.  In RS, the ledger provides a single
observer-independent record: all observers read the same
committed entries, and no contradiction arises.
\end{remark}

%% ============================================================
\section{The Quantum-to-Classical Transition}
\label{sec:decoherence}
%% ============================================================

\begin{proposition}[Decoherence from $J$-Cost Scaling --- Order-of-Magnitude Argument]
\label{prop:decoherence}
Macroscopic superpositions decohere on timescales vastly shorter
than any observable timescale, because the $J$-cost of maintaining
entanglement scales as $N^2$ with particle number.
\end{proposition}

\begin{proof}[Structural argument (order of magnitude)]
Consider a system of $N$ particles.  A \emph{product state}
has total $J$-cost $J_{\mathrm{prod}} \sim N$ (sum of $N$ independent
terms of order 1).  A fully \emph{entangled state} (all $N$ entries
correlated) incurs off-diagonal correlation costs from all $\binom{N}{2}
\sim N^2/2$ pairs, giving $J_{\mathrm{ent}} \sim N^2$.  The cost
differential $\Delta J = J_{\mathrm{ent}} - J_{\mathrm{prod}} \sim
N(N-1)$ drives decoherence: the system minimizes $J$-cost by
transitioning from entangled to product states at a rate
proportional to $\Delta J$.  For $N \sim 10^{23}$ (macroscopic),
$\Delta J \sim 10^{46}$, and decoherence is effectively instantaneous.
\end{proof}

\begin{remark}[Status of the $N^2$ scaling]
The $J$-cost $\sim N^2$ scaling for the entangled state counts
the $\binom{N}{2}$ pairwise correlation terms.  This is a
\emph{structural estimate}; a rigorous derivation of the decoherence
rate from the full RS Hamiltonian (with precise coupling constants)
is an identified open problem.  The parametric argument correctly
captures the essential physics: entanglement maintenance becomes
increasingly costly at large $N$, providing a natural $J$-cost-based
explanation of the quantum-to-classical transition without a
Heisenberg cut.
\end{remark}

\begin{remark}
This explains why quantum effects are observable at microscopic
scales ($N \sim 1$, $\Delta J \sim 0$) but not at macroscopic
scales ($N \gg 1$, $\Delta J \gg 1$).  The transition is smooth,
not sharp---there is no ``Heisenberg cut.''
\end{remark}

%% ============================================================
\section{Comparison with Interpretations}
\label{sec:comparison}
%% ============================================================

\begin{center}
\renewcommand{\arraystretch}{1.3}
\begin{tabular}{@{}lccccc@{}}
\toprule
& \textbf{Copen.} & \textbf{MWI} & \textbf{Bohm} &
\textbf{GRW} & \textbf{RS} \\
\midrule
Collapse & Observer & None & Epistemic &
Random & Commit \\
$\psi$ status & Epistemic & Ontic & Ontic &
Ontic & Inter-tick \\
Born rule & Postulated & Derived* & Equivariant &
Derived & Derived \\
Definite & Observer & All occur & Particles &
Stochastic & Discrete \\
Evolution reversible & Yes & Yes & Yes &
No & No \\
Free params & 0 & 0 & 0 & 2 & 0 \\
\bottomrule
\end{tabular}
\end{center}

\noindent *MWI derivations of the Born rule remain
contested~\cite{Kent,Zurek}.

%% ============================================================
\section{Predictions and Falsifiers}
\label{sec:falsifiers}
%% ============================================================

\begin{enumerate}
  \item \textbf{No macroscopic superpositions}: Detection of a
  macroscopic superposition persisting beyond $8\tau_0/N$ (where
  $N$ is the particle count) would challenge the commit model.

  \item \textbf{Born rule exact}: Any violation of
  $P = |\psi|^2$ at any scale would falsify the $J$-cost
  derivation.

  \item \textbf{Measurement time}: The minimum measurement duration
  is $8\tau_0$.  If measurements could be completed in fewer than
  8 ticks, the epoch structure would be falsified.

  \item \textbf{Irreversibility}: The commit is irreversible.
  Any demonstration of ``uncollapse''---reversing a definite
  measurement outcome to a superposition---would falsify RS.

  \item \textbf{No Heisenberg cut}: The quantum-to-classical
  transition is smooth (governed by $N^2$ scaling), not sharp.
  Discovery of a sharp boundary would challenge the $J$-cost
  decoherence mechanism.
\end{enumerate}

%% ============================================================
\section{Conclusion}
\label{sec:conclusion}
%% ============================================================

The measurement problem is dissolved: recognition \emph{is}
measurement.  The ledger commits to definite outcomes at each
8-tick epoch boundary.  The Born rule is the unique probability
assignment satisfying five structural constraints: non-negativity,
normalization, $J$-cost phase invariance, additivity under
coarse-graining, and continuity.  Definite outcomes are forced by ledger
discreteness.  Collapse is the irreversible commit operation.
Schr\"odinger's cat is never in a superposition.  Wigner's friend
reads the same ledger.  The quantum-to-classical transition is
governed by $J$-cost scaling ($N^2$ vs.\ $N$).  No observer, no
many worlds, no pilot wave, no spontaneous collapse---just the
ledger.

%% ============================================================
\begin{thebibliography}{99}

\bibitem{Aczel1966}
J.~Acz\'el, \emph{Lectures on Functional Equations and Their Applications},
Academic Press, New York (1966).

\bibitem{RS_Axioms}
J.~Washburn, ``The Algebra of Reality: A Recognition Science Derivation
of Physical Law,'' \emph{Axioms} \textbf{15}(2), 90 (2026).
\texttt{doi:10.3390/axioms15020090}

\bibitem{vN}
J.~von Neumann, \emph{Mathematical Foundations of Quantum
Mechanics} (Princeton University Press, 1932).

\bibitem{Bell}
J.~S.~Bell, \emph{Speakable and Unspeakable in Quantum Mechanics}
(Cambridge University Press, 1987).

\bibitem{Everett}
H.~Everett, ``Relative State Formulation of Quantum Mechanics,''
Rev.\ Mod.\ Phys.\ \textbf{29}, 454 (1957).

\bibitem{Bohm}
D.~Bohm, ``A Suggested Interpretation of the Quantum Theory in
Terms of `Hidden' Variables,'' Phys.\ Rev.\ \textbf{85}, 166
(1952).

\bibitem{GRW}
G.~C.~Ghirardi, A.~Rimini, and T.~Weber, ``Unified Dynamics for
Microscopic and Macroscopic Systems,'' Phys.\ Rev.\ D \textbf{34},
470 (1986).

\bibitem{Fuchs}
C.~A.~Fuchs, ``QBism, the Perimeter of Quantum Bayesianism,''
arXiv:1003.5209 (2010).

\bibitem{Kent}
A.~Kent, ``One World versus Many: The Inadequacy of Everettian
Accounts of Evolution, Probability, and Scientific Confirmation,''
in \emph{Many Worlds?}, Oxford University Press (2010).

\bibitem{Zurek}
W.~H.~Zurek, ``Decoherence, Einselection, and the Quantum
Origins of the Classical,'' Rev.\ Mod.\ Phys.\ \textbf{75}, 715
(2003).

\bibitem{FrauchigerRenner}
D.~Frauchiger and R.~Renner, ``Quantum Theory Cannot Consistently
Describe the Use of Itself,'' Nature Comm.\ \textbf{9}, 3711 (2018).

\end{thebibliography}

\end{document}
